% ============================================================
%  Manual de Usuario — Churn Dashboard (Olist)
% ============================================================
\documentclass[11pt]{article}

\usepackage[margin=1in]{geometry}
\usepackage[T1]{fontenc}
\usepackage[utf8]{inputenc}
\usepackage{lmodern}
\usepackage{microtype}
\usepackage{graphicx}
\usepackage{float}
\usepackage{caption}
\usepackage{subcaption}
\usepackage{booktabs}
\usepackage{array}
\usepackage{enumitem}
\usepackage{xcolor}
\usepackage{hyperref}
\hypersetup{colorlinks=true,linkcolor=blue,urlcolor=blue}
\usepackage{bookmark}

\definecolor{lowgreen}{RGB}{22,163,74}
\definecolor{medorange}{RGB}{234,88,12}
\definecolor{highred}{RGB}{220,38,38}

\title{Manual de Usuario\\
\large Churn Dashboard --- Predicción de Abandono de Clientes Olist}
\author{Equipo Microproyecto}
\date{Marzo 2026}

\begin{document}
\maketitle
\tableofcontents
\newpage

% ============================================================
\section{Introducción}

El \textbf{Churn Dashboard} es una herramienta web que permite estimar la
probabilidad de que un cliente de Olist \textbf{no vuelva a comprar} en los
próximos 90 días, a partir de su historial de transacciones. La predicción
es generada por un modelo de Regresión Logística entrenado con datos reales
del marketplace.

Este manual está dirigido a usuarios finales (analistas, equipos de retención)
que deseen usar el tablero sin necesidad de acceder al código fuente.

\subsection*{Acceso al tablero}

El tablero está desplegado en una instancia AWS EC2. Para acceder, abra su
navegador web (Chrome, Firefox o Edge) e ingrese la URL provista por el
administrador, que tiene el formato:

\begin{center}
  \texttt{http://<IP-PUBLICA-EC2>:3000}
\end{center}

No se requiere instalación local ni credenciales para usar el tablero.

% ============================================================
\section{Interfaz principal}

Al cargar el tablero verá dos áreas:

\begin{enumerate}[itemsep=4pt]
  \item \textbf{Panel lateral izquierdo} (\textit{Customer Data}): formulario
        de entrada de datos del cliente.
  \item \textbf{Panel central} (\textit{Customer Churn Prediction Dashboard}):
        resultados de la predicción y análisis de factores.
\end{enumerate}

\begin{figure}[H]
  \centering
  \includegraphics[width=0.95\textwidth]{evidence/dashboard/EVID_DASH_01_prediction.png}
  \caption{Vista principal del tablero con predicción activa.}
\end{figure}

% ============================================================
\section{Panel lateral: ingreso de datos del cliente}

El panel izquierdo está dividido en tres secciones. Complete todos los campos
antes de solicitar la predicción.

\subsection{RFM Analysis}

\begin{center}
\begin{tabular}{lp{9cm}}
\toprule
\textbf{Campo} & \textbf{Descripción} \\
\midrule
\texttt{Recency} &
  Número de días desde la \textbf{última compra} del cliente hasta la fecha
  de análisis. Valores altos indican mayor inactividad. \\[4pt]
\texttt{Frequency} &
  \textbf{Total de órdenes} realizadas históricamente por el cliente. Mínimo: 1. \\[4pt]
\texttt{Monetary} &
  \textbf{Gasto acumulado} del cliente en todas sus órdenes (en la moneda del
  dataset, BRL). \\
\bottomrule
\end{tabular}
\end{center}

\subsection{Order Metrics}

\begin{center}
\begin{tabular}{lp{9cm}}
\toprule
\textbf{Campo} & \textbf{Descripción} \\
\midrule
\texttt{Avg. Items per Order} &
  Promedio de artículos por orden. \\[4pt]
\texttt{Avg. Price Sum} &
  Promedio de la suma de precios por orden (sin flete). \\[4pt]
\texttt{Avg. Freight Sum} &
  Promedio del costo de flete por orden. \\
\bottomrule
\end{tabular}
\end{center}

\subsection{Experience Metrics}

\begin{center}
\begin{tabular}{lp{9cm}}
\toprule
\textbf{Campo} & \textbf{Descripción} \\
\midrule
\texttt{Avg. Review Score} &
  Promedio del puntaje de reseña del cliente (escala 1--5). \\[4pt]
\texttt{Avg. Delivery Days} &
  Promedio de días entre la compra y la entrega real. \\[4pt]
\texttt{Avg. Late Days} &
  Promedio de días de retraso respecto a la fecha estimada de entrega.
  Valores negativos indican entregas anticipadas. \\[4pt]
\texttt{Customer State} &
  Estado de Brasil donde reside el cliente. Seleccione de la lista desplegable
  (AC, AL, AM, \ldots, SP, TO). \\
\bottomrule
\end{tabular}
\end{center}

% ============================================================
\section{Obtener una predicción}

\begin{enumerate}[itemsep=4pt]
  \item Complete todos los campos del panel lateral con los datos del cliente.
  \item Haga clic en el botón \textbf{``Predict Churn''} al pie del formulario.
  \item El panel central se actualizará automáticamente con el resultado.
\end{enumerate}

\noindent\textit{Nota: los campos no pueden quedar vacíos. Si un valor es
desconocido, use el promedio histórico del cliente o un valor representativo
de su segmento.}

% ============================================================
\section{Interpretación de resultados}

\subsection{Tarjeta de probabilidad de churn}

El indicador principal muestra la \textbf{probabilidad de churn} como un
porcentaje en una escala radial. El color del indicador refleja el nivel de
riesgo:

\begin{center}
\begin{tabular}{ccp{7cm}}
\toprule
\textbf{Color} & \textbf{Nivel} & \textbf{Interpretación} \\
\midrule
\textcolor{lowgreen}{\textbf{Verde}} & Bajo &
  Probabilidad $< 40\%$. El cliente tiene alta probabilidad de regresar. \\[4pt]
\textcolor{medorange}{\textbf{Naranja}} & Medio &
  Probabilidad entre $40\%$ y $70\%$. Riesgo moderado; considerar acciones
  de retención preventivas. \\[4pt]
\textcolor{highred}{\textbf{Rojo}} & Alto &
  Probabilidad $\geq 70\%$. El cliente tiene alta probabilidad de no volver.
  Priorizar acciones de retención inmediatas. \\
\bottomrule
\end{tabular}
\end{center}

\subsection{Tarjeta de retención}

A la derecha del indicador de churn aparece la \textbf{probabilidad de retención} (complementaria): $p(retencion) = 1 - p(churn)$.

\subsection{Panel de contribución de features}

\begin{figure}[H]
  \centering
  \includegraphics[width=0.15\textwidth]{evidence/dashboard/EVID_DASH_03_features_export.png}
  \caption{Panel de contribución de features y exportación.}
\end{figure}

Debajo del indicador principal aparece la lista de las \textbf{8 variables
con mayor impacto} en la predicción actual:

\begin{itemize}[itemsep=2pt]
  \item Una barra verde con flecha hacia abajo (\textit{TrendingDown}) indica
        que esa variable \textbf{reduce} la probabilidad de churn
        (favorece la retención).
  \item Una barra roja con flecha hacia arriba (\textit{TrendingUp}) indica
        que esa variable \textbf{aumenta} la probabilidad de churn.
  \item Las barras están ordenadas por magnitud de impacto: la variable en
        la primera posición es la de mayor influencia en esta predicción.
\end{itemize}

Este panel permite al equipo de retención entender \textit{por qué} el modelo
clasifica a ese cliente como en riesgo y tomar acciones concretas
(por ejemplo: si \texttt{Recency} es el factor principal, el cliente lleva
mucho tiempo sin comprar y podría beneficiarse de una oferta de reactivación).

% ============================================================
\section{Visualizaciones macro}

\begin{figure}[H]
  \centering
  \includegraphics[width=0.9\textwidth]{evidence/dashboard/EVID_DASH_02_macro.png}
  \caption{Visualizaciones macro: tasa de churn y distribución por estado.}
\end{figure}

El tablero también incluye visualizaciones a nivel de población:

\begin{itemize}[itemsep=2pt]
  \item \textbf{Distribución de churn}: porcentaje global de clientes con
        churn vs.\ no-churn en el dataset de entrenamiento.
  \item \textbf{Tasa de churn por estado}: mapa o gráfico de barras que
        muestra en qué estados de Brasil la tasa de abandono es más alta,
        útil para segmentación geográfica.
\end{itemize}

% ============================================================
\section{Casos de uso típicos}

\subsection{Evaluar un cliente antes de una campaña de retención}

\begin{enumerate}[itemsep=2pt]
  \item Busque en su CRM los datos del cliente: última compra, número de
        órdenes, gasto total, puntaje de reseñas, etc.
  \item Ingrese los valores en el formulario lateral.
  \item Si el nivel de riesgo es \textbf{Alto} o \textbf{Medio}, incluya
        al cliente en la campaña de retención con prioridad acorde.
  \item Use el panel de contribución para personalizar el mensaje: si
        \texttt{avg\_review\_score} tiene impacto negativo alto, el cliente
        tuvo malas experiencias y puede responder a una disculpa + cupón.
\end{enumerate}

\subsection{Explorar un segmento geográfico}

\begin{enumerate}[itemsep=2pt]
  \item Consulte la visualización macro de tasa de churn por estado.
  \item Seleccione el estado con mayor tasa.
  \item Ingrese valores representativos de clientes de ese estado para
        validar el riesgo individual.
\end{enumerate}

% ============================================================
\section{Preguntas frecuentes}

\begin{description}[style=nextline, itemsep=6pt]
  \item[¿Qué umbral usa el modelo para decidir si hay churn?]
    El modelo fue calibrado con un umbral de $\hat\tau = 0.12$: si la
    probabilidad de recompra supera ese valor, el cliente se clasifica como
    no-churn. En la interfaz, sin embargo, se muestra la probabilidad
    continua para que el analista aplique su propio criterio operativo.

  \item[¿Con qué frecuencia debo actualizar los datos del cliente?]
    Se recomienda usar datos actualizados al momento del análisis. Los
    campos de Recency, Frequency y Monetary son los más sensibles a cambios
    recientes.

  \item[¿Qué hago si la página no carga?]
    Verifique que la instancia EC2 esté activa y que el proceso PM2
    \texttt{churn\_dashboard} esté en estado \texttt{online}. Contacte al
    administrador del sistema.

  \item[¿Puedo exportar el resultado?]
    El resultado de la predicción puede copiarse de la barra de URL o del
    JSON que devuelve el endpoint \texttt{/api/predict}. Consulte al
    administrador para habilitar una exportación estructurada.
\end{description}

\end{document}
